\subsection{Import GPS Trails}
\label{sec:ui_import_gps}

This task allows importing (D)GPS trails into the opened database, by parsing an NMEA file and writing the position updates into the 'RefTraj' DBContent.

\subsubsection{Required NMEA Sentences}

COMPASS uses the NemaTode NMEA parser library, which supports the following NMEA sentence types. All talker ID prefixes are supported (GP, GL, GA, GN --- i.e.\ GPS, GLONASS, Galileo, and multi-GNSS solutions).

\begin{table}[H]
\centering
\begin{tabular}{|l|l|l|}
\hline
\textbf{Sentence} & \textbf{Required} & \textbf{Data used by COMPASS} \\
\hline
GGA & Yes & Time, position (lat/lon), fix quality, number of tracking satellites, altitude \\
\hline
RMC & Yes & Time, \textbf{date}, position, speed over ground, track angle, active/void status \\
\hline
GSA & Yes & Fix type (2D/3D), PDOP, HDOP, VDOP (used for accuracy calculation) \\
\hline
GSV & Optional & Number of visible satellites, satellite SNR (signal-to-noise ratio) \\
\hline
VTG & Optional & Ground speed (used as fallback if RMC speed is unavailable) \\
\hline
\end{tabular}
\caption{NMEA sentences used by COMPASS}
\end{table}

\textbf{Note:} The GGA sentence provides position and time-of-day but \textit{no date}. The date is only available from the RMC sentence. Without RMC, the 'Override Date' option must be used. The GSA sentence provides dilution-of-precision values which are used to compute position accuracy estimates. The GST sentence is \textit{not} supported.

\begin{figure}[H]
  \centering
    \includegraphics[width=16cm,frame]{figures/gps_import_task.png}
  \caption{Import GPS Trail}
\end{figure}

There exist 2 tabs:

\begin{itemize}
\item Main: File path and information text
\item Config: Configuration of data source and secondary information
\end{itemize}
\ \\

\subsubsection{Main Tab}

At the top, a label exists indicating the file to be imported. \\

Below, a text field is given, which after selection of an NMEA file displays the content information and/or error messages. The following information is (commonly) presented:

\begin{itemize}
\item Number of lines parsed
\item Number of (position) fixes
\item Number of skipped fixes, skipped if
\begin{itemize}
\item the time and position is the same as the previous update or
\item the quality is set to 0 (invalid position).
\end{itemize}
\item Number of remaining fixes (to be imported)
\item Number of missing speed vector information
\item Timestamps
\begin{itemize}
\item In NMEA (if certain messages are present) the date information is given
\item The Begin/End timestamps are the first and the last in the file
\item If this information is wrong, it is recommended to either edit the NMEA file, or (if no 24h rollover occurred) to use the 'Override Date' function
\end{itemize}
\end{itemize}

\subsubsection{Config Tab}

\begin{figure}[H]
    \centering
    \includegraphics[width=16cm]{figures/gps_import_config.png}
  \caption{Import GPS Trail Config}
\end{figure}

In the configuration tab, several configuration parameters can be set:

\begin{itemize}
\item SAC: System area code of the reference data source
\item SIC: System identification code of the reference data source
\item Name: Name of the reference data source
\item Time of Day Offset: Time correction factor, in seconds. Set to 0 to disable.
\item Override Date: Set date information, e.g. if missing in NMEA file
\item Mode 3/A Code (octal): Mode 3/A code to be set. Uncheck checkbox to disable.
\item Target Address (hexadecimal): Mode S address to be set. Uncheck checkbox to disable.
\item Callsign: Target identification to be set. Uncheck checkbox to disable.
\item Line ID: Line to be used during import
\end{itemize}
\ \\

\subsubsection{Running}

After selecting an NMEA file, the task can be performed using the 'Import' button.


